% GENERATED FILE. Edit canonical lessons, then run npm run generate:books.
% canonical-source-hash: fnv1a64:c5bcd586e90ec0d4
% canonical-lessons: AR-C158-jadwal, AR-C158-lafita, AR-C158-ijar, AR-C158-ilan, AR-C158-isal, AR-R158-first-pass-looking-back, AR-R158-second-pass-looking-ahead-and-reading-signs

\chapter{A Timetable, a Sign, the Rent, an Advert, a Receipt}
\label{ch:ar-a-timetable-a-sign-the-rent-an-advert-a-receipt}
\hlchaptermodality{158}

\begin{chapteropening}
\textbf{By the end of this chapter:} \emph{I can read a timetable, a sign, an advert and a receipt, and act on what they say.}

Complete the last lesson of chapter 158: I can read a timetable, a sign, an advert and a receipt, and act on what they say.
\end{chapteropening}

\section[\texorpdfstring{\ar{جدول} (jadwal)}{جدول (jadwal)}]{\ar{جدول} (jadwal) — a timetable, a table}

\label{lesson:AR-C158-jadwal}



\begin{quote}
Before the new one: say the Arabic for to decide, then the Arabic for to intend.
\end{quote}

\subsection*{You'll want to know: \ar{جدول}}
\textbf{\ar{جدول}} — \emph{jadwal} — \textquotedblleft{}a timetable, a table\textquotedblright{}.

\emph{Read it:} the board

The board says \textbf{\ar{موعد} \ar{الحافلة}: \ar{الساعة} \ar{الثامنة} \ar{صباحا}} — \emph{mawʿid al-ḥāfila: as-sāʿa ath-thāmina ṣabāḥan} — \textquotedblleft{}Bus time: 8 a.m.\textquotedblright{} Be there before eight.

\begin{cousinweb}
\textbf{\ar{جدول}} (\emph{jadwal}) is built on the native Arabic root \emph{j-d-l}, \textquotedblleft{}a rivulet\textquotedblright{} (Lane). Lane's \emph{jadwal} is a small stream or irrigation channel, and its verb means to rule a book with lines.
\end{cousinweb}

\begin{grammarlens}[title={What you've built}]
One more word for this chapter.
\end{grammarlens}

\subsection*{Your turn}
\begin{itemize}
\raggedright
  \item \emph{Say it:} \emph{jadwal}
  \item \emph{Say it:} \emph{jadwal}, once more
  \item \emph{Say it:} say \emph{jadwal}
  \item \emph{Recall:} say the Arabic for to decide, then the Arabic for to intend, then say \emph{jadwal} again
\end{itemize}

\begin{tcolorbox}[breakable,colback=teal!4,colframe=teal!35!black,title={Before you move on}]
What does \ar{جدول} mean? (\textquotedblleft{}A timetable, a table\textquotedblright{}.) Say it once more.
\end{tcolorbox}
\section[\texorpdfstring{\ar{لافتة} (lāfita)}{لافتة (lāfita)}]{\ar{لافتة} (lāfita) — a sign, a placard}

\label{lesson:AR-C158-lafita}



\begin{quote}
Before the new one: say the Arabic for to intend, then the Arabic for a timetable.
\end{quote}

\subsection*{You'll want to know: \ar{لافتة}}
\textbf{\ar{لافتة}} — \emph{lāfita} — \textquotedblleft{}a sign, a placard\textquotedblright{}.

\emph{Read it:} the sign

The sign says \textbf{\ar{مغلق} \ar{يوم} \ar{الجمعة}} — \emph{mughlaq yawm al-jumʿa} — \textquotedblleft{}Closed on Friday.\textquotedblright{} Come on another day.

\begin{cousinweb}
\textbf{\ar{لافتة}} (\emph{lāfita}) is built on the native Arabic root \emph{l-f-t}, \textquotedblleft{}to turn (the face)\textquotedblright{} (Lane). \emph{Lāfita} is the feminine active participle, \textquotedblleft{}one that turns (the eye)\textquotedblright{}; the placard is a modern sense.
\end{cousinweb}

\begin{grammarlens}[title={What you've built}]
One more word for this chapter.
\end{grammarlens}

\subsection*{Your turn}
\begin{itemize}
\raggedright
  \item \emph{Say it:} \emph{lāfita}
  \item \emph{Say it:} \emph{lāfita}, once more
  \item \emph{Say it:} say \emph{lāfita}
  \item \emph{Recall:} say the Arabic for to intend, then the Arabic for a timetable, then say \emph{lāfita} again
\end{itemize}

\begin{tcolorbox}[breakable,colback=teal!4,colframe=teal!35!black,title={Before you move on}]
What does \ar{لافتة} mean? (\textquotedblleft{}A sign, a placard\textquotedblright{}.) Say it once more.
\end{tcolorbox}
\section[\texorpdfstring{\ar{إيجار} (ījār)}{إيجار (ījār)}]{\ar{إيجار} (ījār) — rent}

\label{lesson:AR-C158-ijar}



\begin{quote}
Before the new one: say the Arabic for a timetable, then the Arabic for a sign.
\end{quote}

\subsection*{You'll want to know: \ar{إيجار}}
\textbf{\ar{إيجار}} — \emph{ījār} — \textquotedblleft{}rent\textquotedblright{}. The money for a flat: \textbf{\ar{إيجار} \ar{الشقة}} — \emph{ījār ash-shaqqa} — \textquotedblleft{}the rent for the flat\textquotedblright{}.

\begin{cousinweb}
\textbf{\ar{إيجار}} (\emph{ījār}) is built on the native Arabic root \emph{ʾ-j-r}, \textquotedblleft{}a reward, a wage; to hire\textquotedblright{} (Lane). The word \emph{ījār} is the verbal noun of Form IV, \textquotedblleft{}letting for hire\textquotedblright{}.
\end{cousinweb}

\begin{grammarlens}[title={What you've built}]
One more word for this chapter.
\end{grammarlens}

\subsection*{Your turn}
\begin{itemize}
\raggedright
  \item \emph{Say it:} \emph{ījār}
  \item \emph{Say it:} \emph{ījār}, once more
  \item \emph{Say it:} say \emph{ījār}
  \item \emph{Recall:} say the Arabic for a timetable, then the Arabic for a sign, then say \emph{ījār} again
\end{itemize}

\begin{tcolorbox}[breakable,colback=teal!4,colframe=teal!35!black,title={Before you move on}]
What does \ar{إيجار} mean? (\textquotedblleft{}Rent\textquotedblright{}.) Say it once more.
\end{tcolorbox}
\section[\texorpdfstring{\ar{إعلان} (iʿlān)}{إعلان (iʿlān)}]{\ar{إعلان} (iʿlān) — an advert, an announcement}

\label{lesson:AR-C158-ilan}



\begin{quote}
Before the new one: say the Arabic for a sign, then the Arabic for rent.
\end{quote}

\subsection*{You'll want to know: \ar{إعلان}}
\textbf{\ar{إعلان}} — \emph{iʿlān} — \textquotedblleft{}an advert, an announcement\textquotedblright{}.

\emph{Read it:} the advert

The advert says \textbf{\ar{شقة} \ar{للإيجار}} — \emph{shaqqa li-l-ījār} — \textquotedblleft{}Flat for rent.\textquotedblright{}

\begin{cousinweb}
\textbf{\ar{إعلان}} (\emph{iʿlān}) is built on the native Arabic root \emph{ʿ-l-n}, \textquotedblleft{}to be public, to be open\textquotedblright{} (Lane). \emph{Iʿlān} is the verbal noun of Form IV, \textquotedblleft{}making public\textquotedblright{}.
\end{cousinweb}

\begin{grammarlens}[title={What you've built}]
One more word for this chapter.
\end{grammarlens}

\subsection*{Your turn}
\begin{itemize}
\raggedright
  \item \emph{Say it:} \emph{iʿlān}
  \item \emph{Say it:} \emph{iʿlān}, once more
  \item \emph{Say it:} say \emph{iʿlān}
  \item \emph{Recall:} say the Arabic for a sign, then the Arabic for rent, then say \emph{iʿlān} again
\end{itemize}

\begin{tcolorbox}[breakable,colback=teal!4,colframe=teal!35!black,title={Before you move on}]
What does \ar{إعلان} mean? (\textquotedblleft{}An advert, an announcement\textquotedblright{}.) Say it once more.
\end{tcolorbox}
\section[\texorpdfstring{\ar{إيصال} (īṣāl)}{إيصال (īṣāl)}]{\ar{إيصال} (īṣāl) — a receipt}

\label{lesson:AR-C158-isal}



\begin{quote}
Before the new one: say the Arabic for rent, then the Arabic for an advert.
\end{quote}

\subsection*{You'll want to know: \ar{إيصال}}
\textbf{\ar{إيصال}} — \emph{īṣāl} — \textquotedblleft{}a receipt\textquotedblright{}. Ask for it when you pay: \textbf{\ar{الإيصال} \ar{من} \ar{فضلك}} — \emph{al-īṣāl min faḍlik} — \textquotedblleft{}The receipt, please.\textquotedblright{}

\begin{cousinweb}
\textbf{\ar{إيصال}} (\emph{īṣāl}) is built on the native Arabic root \emph{w-ṣ-l}, \textquotedblleft{}to arrive, to join\textquotedblright{} (Lane). The word \emph{īṣāl} is the verbal noun of Form IV, \textquotedblleft{}delivering\textquotedblright{}; the receipt is a modern sense.
\end{cousinweb}

\begin{grammarlens}[title={What you've built}]
One more word for this chapter.
\end{grammarlens}

\subsection*{Your turn}
\begin{itemize}
\raggedright
  \item \emph{Say it:} \emph{īṣāl}
  \item \emph{Say it:} \emph{īṣāl}, once more
  \item \emph{Say it:} say \emph{īṣāl}
  \item \emph{Recall:} say the Arabic for rent, then the Arabic for an advert, then say \emph{īṣāl} again
\end{itemize}

\begin{tcolorbox}[breakable,colback=teal!4,colframe=teal!35!black,title={Before you move on}]
What does \ar{إيصال} mean? (\textquotedblleft{}A receipt\textquotedblright{}.) Say it once more.
\end{tcolorbox}
\section[(dialogue)]{Looking back — a first pass}

\label{lesson:AR-R158-first-pass-looking-back}



\begin{quote}
Say the words for an advert and a receipt.
\end{quote}

\subsection*{Your turn}
\begin{itemize}
\raggedright
  \item \emph{Say it:} recently, already; indeed, last, once, to spend
\end{itemize}

\begin{tcolorbox}[breakable,colback=teal!4,colframe=teal!35!black,title={Before you move on}]
Recently? (\textbf{\ar{حديثا}}, \emph{ḥadīthan}.) Already; indeed? (\textbf{\ar{بالفعل}}, \emph{bi-l-fiʿl}.) Last? (\textbf{\ar{الماضي}}, \emph{al-māḍī}.) Once? (\textbf{\ar{ذات} \ar{مرة}}, \emph{dhāta marra}.) To spend? (\textbf{\ar{قضى}}, \emph{qaḍā}.)
\end{tcolorbox}
\section[(dialogue)]{Looking ahead and reading signs — a second pass}

\label{lesson:AR-R158-second-pass-looking-ahead-and-reading-signs}



\begin{quote}
Say the words for next and from now on.
\end{quote}

\subsection*{Your turn}
Say these aloud:

\begin{itemize}
\raggedright
  \item next, from now on, a programme, to decide, to intend
  \item a timetable, a sign, rent, an advert, a receipt
\end{itemize}

\begin{tcolorbox}[breakable,colback=teal!4,colframe=teal!35!black,title={Before you move on}]
Next? (\textbf{\ar{القادم}}, \emph{al-qādim}.) From now on? (\textbf{\ar{من} \ar{الآن} \ar{فصاعدا}}, \emph{min al-ān fa-ṣāʿidan}.) A programme? (\textbf{\ar{برنامج}}, \emph{barnāmaj}.) To decide? (\textbf{\ar{قرر}}, \emph{qarrara}.) To intend? (\textbf{\ar{نوى}}, \emph{nawā}.) A timetable? (\textbf{\ar{جدول}}, \emph{jadwal}.) A sign? (\textbf{\ar{لافتة}}, \emph{lāfita}.) Rent? (\textbf{\ar{إيجار}}, \emph{ījār}.) An advert? (\textbf{\ar{إعلان}}, \emph{iʿlān}.) A receipt? (\textbf{\ar{إيصال}}, \emph{īṣāl}.)
\end{tcolorbox}
